\chapter*{Introducción}
\addcontentsline{toc}{chapter}{Introducción}
\label{ch:introduccion}

% Decisiones ya tomadas y aplicadas a todo el documento:
%
%  - Títulos de los enunciados: se escriben como argumento opcional,
%    \begin{Teorema}[Título]. Los estilos asfplain / asfdefinicion /
%    asfobservacion definidos en main.tex los componen entre paréntesis y
%    abren línea nueva automáticamente, de modo que no hay que poner ni los
%    paréntesis ni el salto a mano.
%
%  - Referencias cruzadas: \ref para teoremas, proposiciones, definiciones,
%    lemas y secciones; \eqref queda reservado para ecuaciones, que es lo
%    único que debe salir entre paréntesis.
%
%  - Notación del semigrupo: T(t) sólo en el capítulo de semigrupos, donde es
%    un semigrupo genérico. A partir del capítulo del problema de Cauchy se
%    escribe e^{At}, por ser la solución de un problema de evolución concreto.

\textcolor{red}{COSAS QUE HACER:}

\textcolor{red}{- FORMATO!!!}

\textcolor{red}{- PONER LEMA, TEOREMA, PROPOSICIÓN, etc CON LA PRIMERA LETRA EN MAYÚSCULA.}

\textcolor{red}{- PONER := EN LAS DEFINICIONES CUANDO TOQUE.}

\textcolor{red}{- PONER EN MAYUSCULA CADA LETRA DEL TITULO DE LOS TEOREMAS, PROPOSICIONES, ETC (O NO, PERO CASARSE CON UNA DECISION). TAMBIEN PENSAR SI PONER Teorema O teorema CUANDO LO REFERENCIO... PENSAR ESTO TAMBIEN PARA LOS TITULOS DE LAS COSAS...}

\textcolor{red}{- PONER TODAS LAS SUCESIONES COMO $( \cdot )$ EN VEZ DE COMO $\{ \cdot \}$.}

\textcolor{red}{- HACER INTROS DE CAPITULOS.}

\textcolor{red}{- ASEGURARME QUE EN LAS DEFINICIONES SOLO HAYA UNA DEFINICION, NO VARIAS.}

\textcolor{red}{- CAMBIAR $||$ POR lVert EN EL CAPÍTULO DE SEMIGRUPOS (EL RESTO YA ESTÁ).}

\section*{Motivación}
\addcontentsline{toc}{section}{Motivación}


\section*{Objetivos}
\addcontentsline{toc}{section}{Objetivos}


\section*{Estructura de la memoria}
\addcontentsline{toc}{section}{Estructura de la memoria}
